\ticket{17}{Ассоциативные правила: Apriori и FP-tree}

\begin{examplan}
    \item Определить ассоциативные правила, поддержку, достоверность и lift.
    \item Описать два этапа задачи: поиск частых наборов и генерация правил.
    \item Объяснить алгоритм Apriori.
    \item Объяснить FP-tree и FP-Growth.
\end{examplan}

\subsection*{Короткая версия для письменного ответа}

\textbf{Обучение ассоциативным правилам} ищет зависимости вида $X\Rightarrow Y$ в транзакционных данных. Классический пример: если покупатель купил хлеб и молоко, то он часто покупает масло.

\subsection*{Основные понятия}

\begin{description}
    \item[Itemset] набор элементов.
    \item[Транзакция] набор элементов, встретившихся вместе.
    \item[Правило] выражение $X\Rightarrow Y$, где $X\cap Y=\varnothing$.
\end{description}

\textbf{Поддержка} показывает, как часто набор встречается:
\[
supp(X\Rightarrow Y)=supp(X\cup Y)=\frac{|\{t\in T\mid X\cup Y\subseteq t\}|}{|T|}.
\]

\textbf{Достоверность} показывает условную вероятность следствия при наличии посылки:
\[
conf(X\Rightarrow Y)=\frac{supp(X\cup Y)}{supp(X)}.
\]

\textbf{Lift} показывает, не объясняется ли правило простой популярностью следствия:
\[
lift(X\Rightarrow Y)=\frac{conf(X\Rightarrow Y)}{supp(Y)}.
\]
Если $lift>1$, между $X$ и $Y$ есть положительная ассоциация.

\subsection*{Общая задача}

\begin{enumerate}
    \item Найти все \textbf{частые наборы}: itemsets с поддержкой не ниже \texttt{minSupport}.
    \item Из частых наборов построить правила с достоверностью не ниже \texttt{minConfidence}.
\end{enumerate}

Первый этап обычно самый дорогой вычислительно.

\subsection*{Apriori}

\textbf{Apriori} основан на антимонотонности поддержки: если набор частый, то все его подмножества тоже частые; если набор нечастый, то все его надмножества нечастые.

Алгоритм:
\begin{enumerate}
    \item посчитать поддержку всех 1-наборов и получить $L_1$;
    \item из $L_{k-1}$ сгенерировать кандидатов $C_k$;
    \item удалить кандидатов, у которых есть нечастое $(k-1)$-подмножество;
    \item просканировать базу и посчитать поддержку кандидатов;
    \item оставить частые $k$-наборы $L_k$;
    \item остановиться, когда $L_k$ пусто.
\end{enumerate}

После этого для каждого частого набора $F$ перебираются непустые подмножества $X\subset F$ и проверяется правило $X\Rightarrow F\setminus X$.

Плюсы Apriori: простота и сильное отсечение кандидатов. Минусы: многократное сканирование базы и взрыв числа кандидатов при низкой поддержке.

\subsection*{FP-tree и FP-Growth}

\textbf{FP-Growth} избегает явной генерации кандидатов. Он сжимает транзакционную базу в FP-дерево.

Построение:
\begin{enumerate}
    \item первым проходом посчитать частоты элементов и оставить частые;
    \item отсортировать частые элементы по убыванию поддержки;
    \item вторым проходом вставить очищенные и отсортированные транзакции в FP-tree;
    \item общие префиксы транзакций разделяют одни и те же пути, поэтому дерево компактно;
    \item заголовочная таблица связывает узлы с одинаковыми элементами.
\end{enumerate}

Поиск частых наборов:
\begin{itemize}
    \item для каждого элемента строится условная база шаблонов;
    \item по ней строится условное FP-дерево;
    \item алгоритм рекурсивно извлекает частые наборы;
    \item если дерево содержит один путь, все комбинации элементов пути являются частыми.
\end{itemize}

FP-Growth обычно быстрее Apriori на плотных данных и при низком пороге поддержки, но сложнее в реализации и может требовать много памяти.

\begin{important}
    \item Ассоциативное правило имеет вид $X\Rightarrow Y$.
    \item Support отвечает за частоту, confidence --- за надежность правила, lift --- за неслучайность связи.
    \item Apriori использует антимонотонность поддержки.
    \item FP-Growth строит FP-tree и не генерирует кандидатов явно.
    \item Правила строятся только из частых наборов и затем фильтруются по confidence.
\end{important}

